% 05-hinge.tex -- What we proved, what we didn't, handoff to chapter 03
%
% Scope: summary of empirical findings; what remains unanswered; explicit
% handoff to Pascanu 2013 chapter.

\section{Điều chương này chứng minh, và điều nó chưa}
\label{sec:ch02-hinge}

Chương này không có định lý. Nó có ba bảng số liệu và hai hộp cầu nối. Nhưng
những con số đó nói một điều rất rõ: RNN với hàm kích hoạt $\tanh$ và huấn
luyện bằng BPTT
\textbf{không học được \tn{phụ thuộc dài hạn}{long-term dependencies}} khi
chuỗi vượt quá một độ dài nhất định. Ở
\tn{bài toán sao chép}{copy task} với $T_{\text{mem}} = 20$, gradient
tại bước đầu tiên yếu hơn gradient tại bước cuối cùng 25 lần. Loss sau 2000
epoch huấn luyện cao hơn loss trước khi huấn luyện.

Đó là hiện tượng quan sát được. Còn nguyên nhân thì sao?

Nguyên nhân là \textbf{cấu trúc đại số của BPTT}. Mỗi bước lùi nhân gradient
với $\mat{W}_{hh}\transpose$ và với đạo hàm của $\tanh$. Nếu trị riêng lớn
nhất của $\mat{W}_{hh}$ nhỏ hơn 1, tích của $T$ ma trận như vậy tiến về 0
khi $T$ lớn. Gradient từ bước cuối không đến được bước đầu. Các tham số ở
gần đầu chuỗi không nhận được tín hiệu học có ý nghĩa, và mạng không thể
học cách biểu diễn thông tin từ quá khứ xa.

Nhưng \textbf{tại sao} trị riêng của $\mat{W}_{hh}$ quyết định mọi thứ?
Điều kiện chính xác cho vanishing và exploding gradient là gì? Tại sao
clipping gradient giải quyết được vấn đề bùng nổ nhưng không giải quyết
được vấn đề triệt tiêu? Và quan trọng nhất: có cách nào để huấn luyện RNN
mà không cần thay đổi kiến trúc không?

Đó là những câu hỏi mà chương tiếp theo trả lời. Chương~3 sẽ đọc bài báo
của Pascanu, Mikolov và Bengio~\autocite{pascanu2013difficulty}, dẫn xuất điều kiện
đủ cho vanishing (qua trị riêng lớn nhất của $\mat{W}_{\text{rec}}$) và điều
kiện cần cho exploding, rồi đánh giá hai giải pháp họ đề xuất: clipping
gradient và bộ chính quy hóa bảo toàn norm. Một trong hai giải pháp đó còn
sống đến năm 2026. Cái còn lại thì không.

Nếu bạn đang đọc cuốn sách này một mạch, đây là lúc để dừng và làm bài tập
cuối chương. Chương~3 sẽ bắt đầu bằng việc mở bài báo của Pascanu, và bạn
sẽ cần những gì bạn vừa thấy; những con số, những norm gradient, cảm giác
về việc một RNN thất bại trông như thế nào; để hiểu vì sao bài báo đó
quan trọng.
